% English revision, 10 September 2026.
\begin{frame}{Task 6: filter geometry and learnable pooling}

\vspace{3mm}\textbf{Two extensions of our CellCNN baseline}
\begin{center}
\begin{tikzpicture}[>=Stealth,every node/.style={font=\small,align=center}]
\node (x) at (0,0) {Cell set\\$N\times p$};
\node[text=Teal] (f) at (3,0) {Mahalanobis-ReLU\\or Quadratic-ReLU};
\node[text=Indigo] (p) at (7,0) {Learnable\\Soft Top Fraction};
\node (o) at (10.5,0) {Sample label\\CMV$^+$/CMV$^-$};
\draw[->,thick] (x)--(f);\draw[->,thick] (f)--(p);\draw[->,thick] (p)--(o);
\node[text=Muted,font=\footnotesize] at (3,-1.05) {Which cell profiles respond?};
\node[text=Muted,font=\footnotesize] at (7,-1.05) {How broadly do we average?};
\end{tikzpicture}\end{center}
Our baseline: linear ReLU filters with \textbf{fixed top-1\,\% pooling}.\\

\takeaway{Predict the sample phenotype from an unordered set of cells.}

\source{Arvaniti \& Claassen (2017), Methods; project snapshot, 9 September 2026: 04c, GrHa/06d, GrHa-learnable-pooling/06b.}
\end{frame}

\begin{frame}{Starting point: an isotropic Gaussian filter}

\begin{columns}[T,onlytextwidth]\column{.60\textwidth}
\[g_k(x)=\exp\!\left(-\frac{\|x-c_k\|^2}{2\ell_k^2}\right)\]
A smooth response around a prototype $c_k$, with one width $\ell_k$ for all marker directions.
\begin{itemize}
\item Equal deviations receive equal penalties in standardized coordinates.
\item Standardization equalizes scale, not predictive relevance.
\item One width cannot allow broad variation in one marker and narrow variation in another.
\end{itemize}
\column{.35\textwidth}\centering
\textbf{Equal-response contours}\par
\begin{tikzpicture}[scale=.65,>=Stealth]
\draw[->,Muted] (-2.2,0)--(2.5,0) node[right,font=\scriptsize] {$x_1$};
\draw[->,Muted] (0,-2.2)--(0,2.4) node[above,font=\scriptsize] {$x_2$};
\draw[Teal,thick] (0,0) circle (1.8);\draw[Teal,thick] (0,0) circle (1.15);
\fill[Ink] (0,0) circle (2pt);\node[above right,font=\small] at (0,0) {$c_k$};
\end{tikzpicture}\par\small Spheres in $p$ dimensions.
\end{columns}
\takeaway{Isotropy is a useful baseline. We relax its equal-tolerance assumption with learned marker weights.}

\source{Gaussian radial basis function: Rasmussen \& Williams (2006), GPML, Ch. 4. Model motivation, not an experimental finding.}
\end{frame}

\begin{frame}{Diagonal Mahalanobis distance}

\[\boxed{d_k^2(x)=\frac1p\sum_{j=1}^p a_{kj}(x_j-c_{kj})^2
=\frac1p(x-c_k)^\top A_k(x-c_k)}\]
\[A_k=\operatorname{diag}(a_{k1},\ldots,a_{kp}),\qquad a_{kj}>0,\qquad \frac1p\sum_j a_{kj}=1.\]
\begin{itemize}
\item $c_k$: preferred marker profile, learned through sample labels.
\item $a_{kj}$: sensitivity to deviations along marker axis $j$.
\item Larger $a_{kj}$ means narrower tolerance, not higher expression.
\end{itemize}
\takeaway{Classical Mahalanobis uses $M=\Sigma^{-1}$. Here $M=A/p$ is learned discriminatively, rather than estimated as an inverse covariance.}

\source{Project definition: GrHa-learnable-pooling/06b. Positive parameterization and normalization: backup slide.}
\end{frame}

\begin{frame}{Diagonal or full Mahalanobis?}

\begin{columns}[T,onlytextwidth]\column{.48\textwidth}\centering
\textcolor{Teal}{\textbf{Diagonal: axis-aligned}}\par
\begin{tikzpicture}[scale=.68,>=Stealth]
\draw[->,Muted] (-2.6,0)--(2.8,0) node[right,font=\scriptsize] {$x_1$};
\draw[->,Muted] (0,-1.9)--(0,2) node[above,font=\scriptsize] {$x_2$};
\draw[Teal,very thick,fill=Teal!8] (0,0) ellipse (2.1 and .9);
\fill[Ink] (0,0) circle(2pt);\node[above right,font=\scriptsize] at (0,0) {$c$};
\end{tikzpicture}\par\small $a_1\delta_1^2+a_2\delta_2^2=\text{constant}$
\column{.48\textwidth}\centering
\textcolor{Indigo}{\textbf{Full: rotation is possible}}\par
\begin{tikzpicture}[scale=.68,>=Stealth]
\draw[->,Muted] (-2.6,0)--(2.8,0) node[right,font=\scriptsize] {$x_1$};
\draw[->,Muted] (0,-1.9)--(0,2) node[above,font=\scriptsize] {$x_2$};
\draw[Indigo,very thick,fill=Indigo!8,rotate=32] (0,0) ellipse (2.1 and .9);
\fill[Ink] (0,0) circle(2pt);\node[above right,font=\scriptsize] at (0,0) {$c$};
\end{tikzpicture}\par\small $a_1\delta_1^2+2a_{12}\delta_1\delta_2+a_2\delta_2^2=\text{constant}$
\end{columns}
\small At $p=37$: \textbf{37} diagonal versus \textbf{703} independent symmetric matrix entries per filter, before normalization; plus 37 center coordinates.
\takeaway{Markers can be coupled, but learning all cross-terms from only 20 donors risks fitting donor-specific geometry.}

\source{Own schematic contours; $703=37\cdot38/2$. Positive definite quadratic forms: Stanford EE263, Lecture 16.}
\end{frame}

\begin{frame}{Cell responses compared}

\renewcommand{\arraystretch}{1.8}
\begin{tabular}{@{}p{.20\linewidth}p{.45\linewidth}p{.29\linewidth}@{}}\toprule
Filter & Cell response $h_k(x)$ & Preferred structure\\\midrule
Linear & $\relu(b_k+w_k^\top x)$ & A direction in marker space\\
Mahalanobis & $\relu\!\left(\rho_k-\frac1p(x-c_k)^\top A_k(x-c_k)\right)$ & Proximity to a center\\
Quadratic & $\relu\!\left(b_k+w_k^\top x+\sum_j q_{kj}x_j^2\right)$ & Free curvature along each axis\\\bottomrule
\end{tabular}
\vspace{3mm}
Each filter applies the same parameters to every cell.\\
Quadratic uses $Q_k=\operatorname{diag}(q_{k1},\ldots,q_{kp})$: no cross-terms, no rotated quadratic axes.
\takeaway{The preactivation defines the geometry; ReLU alone does not create a bounded cell population.}

\source{Project: GrHa/04c, 06b and 06d. Both target variants pool ReLU responses with Soft Top Fraction.}
\end{frame}

\begin{frame}{Mahalanobis is a restricted quadratic form}

\textbf{Expand the preactivation:}
\[\rho-\frac1p(x-c)^\top A(x-c)
=\underbrace{-\frac1p x^\top Ax}_{\text{quadratic}}
+\underbrace{\frac2p c^\top Ax}_{\text{linear}}
+\underbrace{\rho-\frac1p c^\top Ac}_{\text{constant}}.\]
\begin{columns}[T,onlytextwidth]\column{.48\textwidth}
\textcolor{Teal}{\textbf{Mahalanobis-ReLU}}\par
$q_j=-a_j/p<0$, $w_j=2a_jc_j/p$.\par
Negative curvature and normalized scale impose a bounded region around $c$.
\column{.48\textwidth}
\textcolor{Indigo}{\textbf{Quadratic-ReLU}}\par
$q_j$, $w_j$ and $b$ are free parameters.\par
Open regions are possible; $q=0$ recovers the linear filter.
\end{columns}
\takeaway{The quadratic filter contains this Mahalanobis preactivation as a special case, without enforcing its geometry.}

\source{Own algebraic derivation; implementation: GrHa/06d. Here a filter is a cell-scoring function, not a pairwise Mercer kernel.}
\end{frame}

\begin{frame}{Which shapes can a quadratic filter detect?}

\begin{columns}[T,onlytextwidth]\column{.32\textwidth}\centering
\textbf{Linear}\par
\begin{tikzpicture}[scale=.57]
\fill[Indigo!12] (.2,-1.8) rectangle (2,1.8);
\draw[Muted,->] (-2,0)--(2.2,0);\draw[Muted,->] (0,-1.8)--(0,2);
\draw[Indigo,thick] (.2,-1.8)--(.2,1.8);
\end{tikzpicture}\par\small $z=x_1-0.2$\par Half-space
\column{.32\textwidth}\centering\textbf{Negative curvature}\par
\begin{tikzpicture}[scale=.57]
\draw[Muted,->] (-2,0)--(2.2,0);\draw[Muted,->] (0,-1.8)--(0,2);
\draw[Teal,thick,fill=Teal!12] (0,0) circle (1.45);
\end{tikzpicture}\par\small $z=2.1-x_1^2-x_2^2$\par Inside a circle
\column{.32\textwidth}\centering\textbf{Mixed signs}\par
\begin{tikzpicture}[scale=.57]
\fill[Amber!14] plot[domain=-1.8:1.8,samples=45] ({sqrt(1+\x*\x)},\x)--(2.3,1.8)--(2.3,-1.8)--cycle;
\fill[Amber!14] plot[domain=-1.8:1.8,samples=45] ({-sqrt(1+\x*\x)},\x)--(-2.3,1.8)--(-2.3,-1.8)--cycle;
\draw[Muted,->] (-2.3,0)--(2.4,0);\draw[Muted,->] (0,-1.8)--(0,2);
\draw[Amber,thick] plot[domain=-1.8:1.8,samples=45] ({sqrt(1+\x*\x)},\x);
\draw[Amber,thick] plot[domain=-1.8:1.8,samples=45] ({-sqrt(1+\x*\x)},\x);
\end{tikzpicture}\par\small $z=x_1^2-x_2^2-1$\par Open regions
\end{columns}
\vspace{3mm}
Shading: $z(x)>0$, hence positive ReLU response.\\
Without cross-terms $x_ix_j$, the quadratic axes cannot rotate freely.
\takeaway{Quadratic can detect a circular region. Greater flexibility does not guarantee better prediction.}

\source{Own analytical examples, not cell measurements. Boundaries follow from the displayed polynomials.}
\end{frame}




\begin{frame}{Hard top-$k$ pooling}
\textbf{Keep the $k$ largest cell responses and calculate their mean.}
\[s_{(1)}\ge s_{(2)}\ge\cdots\ge s_{(N)},\qquad
\boxed{P_{\mathrm{top}\text{-}k}=\frac1k\sum_{r=1}^{k}s_{(r)}}\]
\vspace{3mm}
For a top fraction $\alpha$, choose
\[k=\max(1,\lfloor\alpha N\rfloor).
\qquad\text{Our baseline: }\alpha=0.01.\]
\vspace{3mm}
\textbf{Problem: the selection size is discrete.}\\
Changing $\alpha$ either leaves $k$ unchanged or adds a whole cell.
\takeaway{The hard selection is not differentiable at changes in $k$; between them, its gradient with respect to $\alpha$ is zero.}
\source{Project top-$k$ mean pooling: GrHa/04c and 06d. For fixed $k$, selected response values still receive gradients.}
\end{frame}


\begin{frame}{Relax the hard selection}
\textbf{Hard top-$k$ can be written as a binary mask.}
\[m_r^{\mathrm{hard}}=\begin{cases}1,&r\le k\\0,&r>k\end{cases}
\qquad\text{Keep the first $k$ sorted responses.}\]
\vspace{3mm}
\textbf{Replace the abrupt decision with a smooth transition.}
\[\underbrace{m_r^{\mathrm{hard}}\in\{0,1\}}_{\text{keep or discard}}
\quad\longrightarrow\quad
\underbrace{m_r=\operatorname{sigmoid}\!\left(\frac{\alpha-u_r}{\tau}\right)}_{\text{smooth rank mask}}\]
\small $u_r$: relative rank position; $\alpha$: fraction cutoff; $\tau$: transition width.
\takeaway{The sigmoid gives high-response ranks values near one and lower ranks values near zero.}
\source{Smooth relaxation of rank-based selection. The complete pooling operator is defined on the next slide.}
\end{frame}

\begin{frame}{Soft Top Fraction: definition}
\textbf{For each filter, sort the cell responses:}
\[s_{(1)}\ge\cdots\ge s_{(N)}.\]
\[\underbrace{u_r=\frac{r-\tfrac12}{N}}_{\text{relative rank position}},\qquad
\underbrace{\alpha=\operatorname{sigmoid}(\theta)}_{\text{learned fraction}}\]
\[m_r=\operatorname{sigmoid}\!\left(\frac{\alpha-u_r}{\tau}\right),\qquad
\boxed{P_{\mathrm{soft}}=\frac{\sum_{r=1}^{N}m_rs_{(r)}}{\sum_{r=1}^{N}m_r}}\]
\small Dividing by the mask sum gives a weighted mean. One $\alpha$ is learned per filter and shared across samples; $\tau$ stays fixed.
\takeaway{The weights are differentiable in $\alpha$, allowing the classification loss to learn the pooling fraction.}
\source{GrHa-learnable-pooling/06b: initial $\alpha=0.01$, fixed $\tau=0.002$. Sorting and ReLU are differentiable almost everywhere; derivation in backup.}
\end{frame}

\begin{frame}{SKIP: Same array: hard top-$k$ versus sigmoid selection}
\[ [0,3,0,8,1,6,0,2]\quad\xrightarrow{\text{sort}}\quad [8,6,3,2,1,0,0,0]\]
\textbf{Select the top 25\,\%: $k=2$ for eight cells.}
\begin{center}\renewcommand{\arraystretch}{1.4}
\begin{tabular}{@{}lrrrrrrrr@{}}\toprule
Sorted responses & 8 & 6 & 3 & 2 & 1 & 0 & 0 & 0\\\midrule
Hard selection & \textcolor{Teal}{1} & \textcolor{Teal}{1} & 0 & 0 & 0 & 0 & 0 & 0\\
Sigmoid selection & \textcolor{Teal}{1.000} & \textcolor{Teal}{.958} & \textcolor{Amber}{.042} & .000 & .000 & .000 & .000 & .000\\\bottomrule
\end{tabular}\end{center}
\[P_{\mathrm{top}\text{-}2}=\frac{8+6}{2}=7
\qquad\qquad P_{\mathrm{soft}}=\frac{\sum_r m_rs_{(r)}}{\sum_r m_r}\approx6.937\]
\takeaway{Both focus on the same strongest cells. The sigmoid replaces the abrupt cutoff with a small transition.}
\source{Illustration: $\alpha=.25$, $\tau=.02$ (model: $.002$). Masks are rounded; the soft mean uses unrounded values.}
\end{frame}


\begin{frame}{Regularization: different reference models}

\begin{columns}[T,onlytextwidth]\column{.48\textwidth}
\textcolor{Indigo}{\textbf{Quadratic}}\par
\[\mathcal L=\mathcal L_{\rm CE}+10^{-4}(\|W\|_F^2+\|Q\|_F^2+\|V\|_F^2)\]
\small $Q=(q_{kj})$: quadratic coefficient table.\par
\vspace{2mm}
Initialization at $q=0$ reproduces the linear baseline. L2 on $Q$ favors little additional curvature.
\column{.48\textwidth}
\textcolor{Teal}{\textbf{Mahalanobis, learnable branch}}\par
\[\mathcal L=\mathcal L_{\rm CE}+10^{-4}\|V\|_F^2
+10^{-3}\operatorname{mean}_{k,j}(a_{kj}-1)^2\]
\small $\operatorname{mean}_j a_{kj}=1$ fixes the scale.\par
\vspace{2mm}
The geometry penalty favors equal axis weights. In the read code, it is applied only to the training loss.
\end{columns}
\vspace{3mm}
\small With normalized $a$, $\operatorname{mean}(a^2)=1+\operatorname{mean}((a-1)^2)$: L2 on normalized weights cannot collapse the distance to zero.
\takeaway{Quadratic is pulled toward a linear filter; Mahalanobis toward isotropic geometry.}

\source{GitHub snapshot, 9 September 2026: GrHa/06d and GrHa-learnable-pooling/06b. Geometry coefficient is provisional; GrHa/06b ReLU + mean uses only output L2.}
\end{frame}

\begin{frame}{How should we score the learned populations?}

\begin{columns}[T,onlytextwidth]\column{.47\textwidth}
\textbf{Cluster geometry}\par
\vspace{2mm}
Davies--Bouldin: compact, separated groups.\par
\vspace{3mm}
A distinct cell type need not distinguish CMV$^+$ from CMV$^-$. 
\column{.47\textwidth}
\textbf{Phenotype association}\par
\vspace{2mm}
Does population frequency distinguish sample labels?\par
\vspace{3mm}
We have labels for samples, not individual cells.
\end{columns}
\vspace{4mm}
\takeaway{We seek predictive, disease-associated cells, rather than a complete cell-type partition. Association does not establish causation.}

\source{Arvaniti \& Claassen (2017); Davies \& Bouldin (1979). Samples are represented as bags of cells.}
\end{frame}

\begin{frame}{Select a predictive filter, then count responding cells}

\textbf{1. Choose the filter with the strongest positive output contrast}
\[\Delta_k=V_{1k}-V_{0k},\qquad k^*=\arg\max_{k:\Delta_k>0}\Delta_k.\]
\small This selects by the learned classification weights, not the largest cell response.
\vspace{2mm}\par
\normalsize\textbf{2. Apply a half-maximum response cutoff}
\[t_{k^*}=\tfrac12\max_{x\in\mathrm{Inner\text{-}Train}}h_{k^*}(x),\qquad
f_b=\frac1{N_b}\sum_i\mathbf1\{h_{k^*}(x_{bi})>t_{k^*}\}.\]
\small Calibrate on training cells; apply the same cutoff to every test donor.
\takeaway{Half-maximum is a heuristic, not a biological boundary. Its suitability must be checked and may require adaptation for our filters.}

\source{Existing selection and frequency rule: GrHa/06d. No positive output contrast: report population analysis as missing. Alternative cutoffs must be chosen without test labels.}
\end{frame}


\appendix
\begin{frame}{Backup: ReLU introduces an absolute activation boundary}

\begin{columns}[T,onlytextwidth]\column{.50\textwidth}
\[h_k(x)=\relu(\rho_k-d_k^2(x))\]
\[\rho_k>0\qquad\text{learned positive threshold}\]
\small Inside $d_k^2<\rho_k$: positive response.\\
Outside: zero response and zero local gradient through ReLU.\par
\vspace{3mm}
$\rho_k$ is learned in \emph{squared-distance units}. The corresponding metric radius is $\sqrt{\rho_k}$.
\column{.46\textwidth}\centering
\begin{tikzpicture}[x=.9cm,y=.9cm,>=Stealth]
\draw[->,Muted] (-.15,0)--(4.9,0) node[right] {$d^2$};
\draw[->,Muted] (0,-.15)--(0,3.5) node[above] {$h$};
\draw[Teal,very thick] (0,3)--(3,0)--(4.6,0);
\draw[dashed,Muted] (3,0)--(3,3);
\node[below] at (3,0) {$\rho$};\node[left] at (0,3) {$\rho$};
\node[font=\small,Teal] at (1.3,2.8) {active};
\node[font=\small,Muted] at (4,1) {inactive};
\end{tikzpicture}
\end{columns}
\takeaway{ReLU can suppress background responses, but a filter that is inactive on all training cells may receive no useful gradient.}

\source{GrHa/06b: ReLU + mean. ReLU + Soft Top Fraction is the target combination}
\end{frame}
\begin{frame}{Backup: Activation, pooling and population thresholds}

\renewcommand{\arraystretch}{1.7}
\begin{tabular}{@{}p{.24\linewidth}p{.33\linewidth}p{.37\linewidth}@{}}\toprule
Quantity & Definition & Role\\\midrule
Activation boundary & $h_k=\relu(\rho_k-d_k^2)$\newline or $h_k=\relu(z_k)$ & Defines positive responses; learned through filter parameters.\\
Rank fraction $\alpha_k$ & $m_{rk}=\sig((\alpha_k-u_r)/\tau)$ & Controls pooling breadth; learned per filter.\\
Population threshold $t_k$ & $t_k=\tfrac12 M_k$\newline $M_k=\max_{\rm Inner\text{-}Train}h_k(x)$ & Heuristic population cutoff; calibrated after training.\\\bottomrule
\end{tabular}
\takeaway{Half-maximum is a separate evaluation heuristic. It is neither the learned rank fraction nor a threshold trained by backpropagation.}

\source{Half-maximum analysis: GrHa/06d and shared discussion. Target Mahalanobis-ReLU model: apply the same rule to $h_k$.}
\end{frame}
\begin{frame}{Backup: Evaluate phenotype association and stability}

\begin{columns}[T,onlytextwidth]\column{.48\textwidth}
\textbf{Association across test donors}
\[\mathrm{AUC}_{\rm frequency}=\mathrm{ROC\text{-}AUC}(y_b,f_b)\]
\[\Delta f=\overline f_{\rm CMV+}-\overline f_{\rm CMV-}\]
\small These complement network AUC. Do not reverse the direction after observing a test AUC below $0.5$.\par
\vspace{2mm}
Each donor is an observation; additional bags do not create independent donors.
\column{.48\textwidth}
\textbf{Useful additional checks}
\begin{itemize}
\item Similar marker profiles across donor splits and random seeds.
\item Actual marker coexpression in high-response cells.
\item Sensitivity to extreme training responses that set the half-maximum threshold.
\end{itemize}
\end{columns}
\takeaway{Frequency AUC measures association; stability and biological plausibility help assess whether the population is reproducible.}

\source{Project metrics: GrHa/06d. Stability and threshold-sensitivity checks are methodological recommendations, not reported results.}
\end{frame}
\begin{frame}{Backup: distance geometry and response}

\begin{columns}[T,onlytextwidth]\column{.48\textwidth}
\textcolor{Teal}{\textbf{Learn the geometry}}\par
\[\|x-c\|^2\quad\longrightarrow\quad (x-c)^\top M(x-c)\]
A diagonal $M$ stretches the contours along marker axes.\par
\vspace{3mm}
With $a_j=1$, our normalized distance reduces to the mean squared Euclidean distance.
\column{.48\textwidth}
\textcolor{Indigo}{\textbf{Choose the response}}\par
\[g(x)=\exp\!\left(-d^2(x)/(2\ell^2)\right)\]
\[h(x)=\relu\!\left(\rho-d^2(x)\right)\]
Gaussian: smooth, positive tails.\\
ReLU: zero response outside a learned boundary.
\end{columns}
\takeaway{Mahalanobis does not require ReLU. We choose ReLU for an explicit inactive region, not because Gaussian responses are inherently worse.}

\source{Own comparison of the stated functions. Gaussian geometry: GPML, Ch. 4; ReLU prototype response: GrHa/06b.}
\end{frame}

\begin{frame}{Backup: a gradient reaches the learned fraction}

\textbf{Within a fixed response ordering}, let $S_k=\sum_r m_{rk}$.
\[\frac{\partial m_{rk}}{\partial\alpha_k}=\frac{m_{rk}(1-m_{rk})}{\tau}\]
\[\boxed{\frac{\partial P_k}{\partial\alpha_k}
=\frac1{\tau S_k}\sum_r m_{rk}(1-m_{rk})(s_{(r)k}-P_k)}\]
\[\frac{\partial\mathcal L}{\partial\theta_k}
=\frac{\partial\mathcal L}{\partial P_k}
\frac{\partial P_k}{\partial\alpha_k}\alpha_k(1-\alpha_k)\]
\small The term $s_{(r)k}-P_k$ accounts for the changing denominator. The loss determines whether the fraction should increase or decrease.
\takeaway{The learning path exists, but gradients may vanish. Sorting and ReLU are differentiable almost everywhere, not everywhere smooth.}

\source{Own quotient-rule derivation. Gradients to responses: $\partial P_k/\partial s_{(r)k}=w_{rk}$ within a fixed ordering.}
\end{frame}

\begin{frame}{Backup: why parameterize the threshold with softplus?}

\[\rho_k=\operatorname{softplus}(\eta_k)+\varepsilon
=\log(1+e^{\eta_k})+\varepsilon>0\]
\begin{itemize}
\item $\eta_k\in\mathbb R$ is an unconstrained parameter updated by the optimizer.
\item Softplus maps it smoothly to a positive squared-distance threshold.
\item $\varepsilon$ keeps the threshold strictly above zero numerically.
\end{itemize}
\[\frac{\partial\rho_k}{\partial\eta_k}=\sig(\eta_k)\]
\takeaway{This is a positivity constraint implemented through a smooth parameterization, not an additional threshold or pooling operation.}

\source{GrHa/06b, radii(). Softplus is distinct from the sigmoid rank mask used in Soft Top Fraction.}
\end{frame}

\begin{frame}{Backup: positive, normalized marker weights}

\begin{columns}[T,onlytextwidth]\column{.48\textwidth}
\textbf{Parameterization}
\[u_{kj}=\operatorname{softplus}(\beta_{kj})+\varepsilon\]
\[a_{kj}=\frac{u_{kj}}{\frac1p\sum_\ell u_{k\ell}}\]
\[a_{kj}>0,\qquad \operatorname{mean}_j a_{kj}=1\]
\column{.48\textwidth}
\textbf{Effect}\par
The common scale of the weights is fixed.\par
\vspace{3mm}
The denominator couples all marker weights.\par
\vspace{3mm}
Normalization is a parameter constraint; the additional geometry penalty favors $a_{kj}\approx1$.
\end{columns}
\takeaway{A diagonal metric still models a joint marker profile, even without explicit cross-terms or rotations.}

\source{GrHa-learnable-pooling/06b, marker\_weights.}
\end{frame}
